
\begin{proof}
\Cref{thm:sheet-cp-protection} gives
\[
\arg\det M_u=\arg\det M_d=0.
\]
By \Cref{lem:topological-sector-positivity}, the admissible partition function admits a sector
decomposition
\[
Z_{\mathrm{adm}}(\theta)=\sum_{Q\in\mathbb Z} Z_Q e^{iQ\theta},
\qquad
Z_Q\ge 0.
\]
The antiunitary sheet symmetry of \Cref{thm:sheet-cp-protection} gives
\[
\mathcal C_\Sigma\,Q\,\mathcal C_\Sigma^{-1}=-Q,
\]
so the sector weights satisfy $Z_Q=Z_{-Q}$. Hence
\[
\begin{aligned}
Z_{\mathrm{adm}}(\theta)
&=Z_0+2\sum_{Q>0}Z_Q\cos(Q\theta)\\
&\le Z_0+2\sum_{Q>0}Z_Q
=Z_{\mathrm{adm}}(0).
\end{aligned}
\]
By \Cref{lem:nontrivial-topological-support}, one has $Z_1>0$. Therefore for every
$\theta\not\equiv 0\mod 2\pi$,
\[
Z_{\mathrm{adm}}(\theta)
\le
Z_{\mathrm{adm}}(0)-2Z_1\bigl(1-\cos\theta\bigr)
<
Z_{\mathrm{adm}}(0),
\]
which implies
\[
F(\theta)>F(0)
\qquad
\text{for every }\theta\not\equiv 0\mod 2\pi.
\]
Thus the unique minimum is at $\theta=0$, and therefore $\bar\theta=0$.
\end{proof}

\begin{theorem}[Reflection positivity on the admissible and geometric sector]
\label{thm:reflection-positivity-admissible}
Under the admissibility complex, the geometric Hodge closure, and the setup of
\Cref{def:relative-generating-functional,lem:admissible-reflection-positivity}, every admissible
or geometric observable $O$ with support in $\tau>0$ satisfies
\[
\langle \Theta O,O\rangle_{\mathrm{rel}}\ge 0.
\]
\end{theorem}

\begin{proof}
Let $O$ be an admissible or geometric observable supported in $\tau>0$. By construction of the
geometric Hodge sector, geometric observables are built from the same admissible positive-time
algebra, possibly with insertions from the BRST/Hodge-reduced geometric block. The hypotheses of
\Cref{lem:admissible-reflection-positivity} therefore apply to every monomial in the field expansion
of $O$, and hence to every finite linear combination of such monomials.

For each such term the reflected quadratic form is nonnegative:
\[
\langle \Theta O,O\rangle_{\mathrm{rel}}\ge 0.
\]
The quotient normalization contributes only the fixed positive denominator
$Z_{\mathrm{ref}}[0,0,0]$, so no sign change can occur. Continuity then extends the same statement
from finite polynomial observables to the completed admissible/geometric observable algebra. Thus
every admissible or geometric observable supported in positive Euclidean time satisfies the claimed
reflection-positivity bound.
\end{proof}

\begin{theorem}[\textnormal{[A]} Lorentzian reconstruction on \texorpdfstring{$\Padm$}{Padm}]
\label{thm:lorentzian-reconstruction}
Under the same hypotheses, the admissible Euclidean theory reconstructs a self-adjoint Lorentzian Hamiltonian on the $\Padm$ sector, including the physical graviton and scalaron blocks of the geometric Hodge sector.
\end{theorem}

\begin{proof}
\Cref{thm:tailored-os-car-reconstruction} constructs the OS/CAR Hilbert space on $\Padm$ together
with its positive-energy Euclidean time semigroup and reconstructed field net. The bosonic graviton
and scalaron blocks are included because the geometric Hodge observables belong to the same
reflection-positive admissible algebra, and the fermionic sector is reconstructed simultaneously by
the CAR part of the same appendix theorem. The generator of the Euclidean time semigroup is
therefore the claimed self-adjoint Lorentzian Hamiltonian on $\Padm$.
\end{proof}

The admissible vacuum question is now treated in four explicit blocks. The local strict positivity
of \Cref{prop:strict-positive-c6-green,cor:uniform-positive-c6-green} is the transport-side input
for the primitive positive-transfer closure, the thermodynamic limit, and the parent-gap
stability chain stated below.

\begin{theorem}[Block 1: Finite-volume admissible vacuum]
\label{thm:admissible-vacuum-existence}
\label{thm:admissible-vacuum-finite-volume}
Assume finite volume or a declared cutoff, weak-* compactness of $\Vadm$, and lower
boundedness plus coercivity of $\Padm\Hrel\Padm$. Then the finite-volume admissible vacuum
exists as
\[
\rho^\ast_V=\argmin_{\rho\in\Vadm}\tr(\rho\Hrel),
\]
and its renormalized seam vacuum energy is denoted by $\Lambda_{\Seam,\mathrm{ren}}^{(V)}$. This
first block establishes existence only; uniqueness is supplied by Block 2.
\end{theorem}

\begin{proof}
The direct method applies on $\Vadm$ under the stated compactness and coercivity assumptions.
The renormalized seam residue supplies the vacuum-energy identification in the same declared
scheme.
\end{proof}

\begin{tcolorbox}[colback=tfptblue!2, colframe=tfptblue!40, breakable,
    title={Primitivity lemma underlying Block 2}, colbacktitle=tfptblue!60]
Let $G_V$ be the directed graph on the retained CP-even admissible basis vectors, with a directed
edge
\[
u\to v
\]
iff the corresponding matrix element of the blocked transfer operator $T_V$ is strictly positive.
Then:
\begin{enumerate}[label=(\roman*), leftmargin=1.5em, topsep=3pt, itemsep=2pt]
    \item by \Cref{prop:strict-positive-c6-green}, every local admissible $C_6$ transport step is
    strictly positive on the retained basis;
    \item by the admissible word grammar and the finite blocked transport grammar, every retained
    basis vector reaches every other one after finitely many blocked steps, so $G_V$ is strongly
    connected;
    \item the diagonal admissibility weights and blocked selector preserve positive return
    amplitudes on retained basis states, so $G_V$ is aperiodic.
\end{enumerate}
Hence there exists $N_V<\infty$ such that
\[
T_V^{N_V}
\]
has strictly positive matrix entries on $\mathcal K_{+,\mathrm{adm}}^{(V)}$. Equivalently, $T_V$
is primitive.
\end{tcolorbox}

\begin{theorem}[Block 2: Uniqueness of the CP-even admissible vacuum from the primitive positive transfer]
\label{thm:cp-even-vacuum-uniqueness}
Fix finite volume or a declared cutoff. Let $\mathcal K_{+,\mathrm{adm}}^{(V)}$ be the cone of
CP-even admissible vectors in the retained basis, and let $T_V$ be the blocked admissible
transfer operator obtained from the local finite $C_6$ kernels
$\mathcal Y_y^{(\varepsilon)}$, the diagonal admissibility weights, and the two-stage selector
$\Padm$. Then $T_V$ is primitive on $\mathcal K_{+,\mathrm{adm}}^{(V)}$. Hence $T_V$ has a unique
strictly positive Perron vector $v_V^\ast$, the finite-volume CP-even admissible vacuum
$\rho^\ast_V$ of \Cref{thm:admissible-vacuum-finite-volume} is unique, and the gap above that
vacuum is strictly positive.
\end{theorem}

\begin{proof}[Proof sketch]
The primitivity lemma stated immediately above turns the transport positivity of
\Cref{prop:strict-positive-c6-green} and the admissible word connectivity into an explicit
graph-theoretic criterion for $T_V$. Perron--Frobenius / Krein--Rutman then gives a unique
strictly positive Perron vector $v_V^\ast$ in the cone
$\mathcal K_{+,\mathrm{adm}}^{(V)}$, which yields the unique CP-even finite-volume vacuum. The
simplicity of the Perron root is precisely the finite-volume gap statement.
\end{proof}

\begin{lemma}[Boundary decoupling for local admissible observables]
\label{lem:boundary-decoupling-admissible}
Let $A$ be a local admissible observable with support contained in the interior of $V$. Assume the
blocked transfer operators along the admissible exhaustion are primitive with a uniform Perron gap
on the orthogonal complement of the Perron line, and let $\epsilon_f>0$ be the transport gap of the
retained branch. Then there exist constants
\[
C_A>0,
\qquad
\mu>0,
\]
independent of $V\subset W$ along the admissible exhaustion, such that
\[
\bigl|\rho_W^\ast(A)-\rho_V^\ast(A)\bigr|
\le
C_A\,e^{-\mu\,\mathrm{dist}(\mathrm{supp}\,A,\partial V)}.
\]
In particular the local vacuum expectations form a Cauchy family on the quasilocal admissible
algebra.
\end{lemma}

\begin{proof}[Proof sketch]
Primitivity gives exponential contraction of the blocked transfer dynamics onto the Perron line,
with rate controlled by the uniform Perron gap. A local observable supported at distance
$\mathrm{dist}(\mathrm{supp}\,A,\partial V)$ from the boundary can therefore detect a change from
$V$ to $W$ only through transfer segments that propagate from the boundary into the support of $A$.
Those segments are suppressed exponentially by the transport gap $\epsilon_f$, while the orthogonal
component is damped exponentially by the Perron contraction. Absorbing the local operator norm and
finite blocking constants into $C_A$ gives the displayed estimate, for instance with any
$\mu<\min\{\epsilon_f,\Delta_{\mathrm{PF}}\}$ once $\Delta_{\mathrm{PF}}$ denotes the uniform
Perron-gap scale.
\end{proof}

\begin{theorem}[Block 3: Thermodynamic limit of the admissible vacuum]
\label{thm:thermodynamic-limit}
Let $(V_n)_{n\ge1}$ be an admissible exhaustion by blocked seam cells, and let $T_{V_n}$ be the
primitive positive blocked transfer operator of \Cref{thm:cp-even-vacuum-uniqueness}. Then the
associated normalized vacuum states converge weak-* to a unique thermodynamic-limit admissible
vacuum $\rho^\ast$ with renormalized seam vacuum energy $\Lambda_{\Seam,\mathrm{ren}}$.
\end{theorem}

\begin{proof}
Fix a local admissible observable $A$ with support contained in a finite region $K$. Choose $n_0$
such that $K\subset \operatorname{int}(V_n)$ for all $n\ge n_0$. By
\Cref{lem:boundary-decoupling-admissible}, for all $m>n\ge n_0$,
\[
\bigl|\rho^*_{V_m}(A)-\rho^*_{V_n}(A)\bigr|
\le
C_A e^{-\mu\,\mathrm{dist}(K,\partial V_n)}.
\]
Since $\mathrm{dist}(K,\partial V_n)\to\infty$ along the exhaustion, the right-hand side tends to
zero. Hence $\rho^*_{V_n}(A)$ is Cauchy in $\mathbb C$ and therefore convergent. Define
\[
\rho^*(A):=\lim_{n\to\infty}\rho^*_{V_n}(A)
\]
for every local admissible observable $A$.

Linearity is immediate. Positivity follows from
\[
\rho^*(A^*A)=\lim_{n\to\infty}\rho^*_{V_n}(A^*A)\ge 0,
\]
and normalization from
\[
\rho^*(1)=\lim_{n\to\infty}\rho^*_{V_n}(1)=1.
\]
Moreover, $|\rho^*(A)|\le \|A\|$ because the same bound holds for each state $\rho^*_{V_n}$. Thus
$\rho^*$ extends uniquely by continuity from the local algebra to the quasilocal admissible
algebra.

To prove uniqueness, let $\widetilde\rho$ be any weak-* cluster point of the family. Then for every
local $A$ one has
\[
\widetilde\rho(A)=\lim_{k\to\infty}\rho^*_{V_{n_k}}(A)=\rho^*(A),
\]
because the local expectations form a Cauchy family with a unique limit. Since local observables are
norm dense in the quasilocal algebra, it follows that $\widetilde\rho=\rho^*$. Hence the weak-*
limit is unique.

The same boundary-decoupling argument applied to the renormalized seam vacuum energy density shows
that the corresponding finite-volume renormalized vacuum energies form a Cauchy family, and
therefore converge to a unique limit denoted $\Lambda_{\Seam,\mathrm{ren}}$.
\end{proof}

\begin{lemma}[Parent-Hamiltonian gap inherited from the primitive transfer block]
\label{lem:parent-gap-from-transfer-block}
Fix finite volume. Let $T_V$ be the primitive blocked transfer operator of
\Cref{thm:cp-even-vacuum-uniqueness} with unique strictly positive Perron vector $v_V^\ast$, and
let $H_{\mathrm{par}}$ be the blocked admissible parent Hamiltonian built from the local
admissibility projectors together with the local Perron vacuum projector onto $v_V^\ast$. Then
$H_{\mathrm{par}}$ has the same unique local vacuum sector and a strictly positive parent gap
\[
m_0:=\operatorname{gap}(H_{\mathrm{par}})>0.
\]
\end{lemma}

\begin{proof}
Every summand of $H_{\mathrm{par}}(V)$ is positive semidefinite by construction. By design, all
local admissibility penalties annihilate the admissible Perron sector, and the Perron vacuum
projector annihilates the distinguished Perron line generated by $v_V^*$. Hence
\[
H_{\mathrm{par}}(V)v_V^*=0,
\]
so $\mathbb C v_V^*$ is contained in the kernel.

Conversely, let $\psi\in\ker H_{\mathrm{par}}(V)$. Then
\[
0=\langle \psi,H_{\mathrm{par}}(V)\psi\rangle
\]
and since $H_{\mathrm{par}}(V)$ is a finite sum of positive semidefinite terms, each summand must
have vanishing expectation on $\psi$. Therefore $\psi$ lies in the kernel of every local penalty
term and in the kernel of the Perron projector term. By the definition of the parent Hamiltonian,
the common zero set of these penalties is exactly the Perron vacuum line. Hence
\[
\ker H_{\mathrm{par}}(V)=\mathbb C v_V^*.
\]

Because the blocked finite-volume Hilbert space is finite dimensional, the spectrum of
$H_{\mathrm{par}}(V)$ is discrete and nonnegative. The zero eigenvalue is simple, so the smallest
nonzero eigenvalue exists and is strictly positive. This number is precisely
\[
m_0(V):=\min\bigl(\operatorname{spec}(H_{\mathrm{par}}(V))\setminus\{0\}\bigr)>0.
\]
\end{proof}

\begin{lemma}[Diagonal admissibility weights separate the retained basis]
\label{lem:diagonal-weights-separate-basis}
Let $\mathcal D_V$ be the commuting family of diagonal admissibility weights entering the blocked
transfer operator on the retained basis of $\mathcal K_{+,\mathrm{adm}}^{(V)}$. If
\[
u\neq v
\]
are two retained basis words, then there exists $D\in\mathcal D_V$ such that
\[
D_{uu}\neq D_{vv}.
\]
\end{lemma}

\begin{proof}
Two distinct retained basis words differ in at least one blocked admissibility datum: local
transport count, carrier label, or retained selector weight. The blocked diagonal admissibility
weights were constructed precisely from these discrete data, so some diagonal weight separates the
two words. Hence the family $\mathcal D_V$ separates the retained basis.
\end{proof}

\begin{theorem}[Injective blocked transfer and uniform parent gap]
\label{thm:uniform-parent-gap-injective}
Let $P_V$ be the blocked positive transport operator on the retained auxiliary space, and let
$\mathcal D_V$ be the family of diagonal admissibility weights. Under
\Cref{prop:strict-positive-c6-green,lem:diagonal-weights-separate-basis}, after finite blocking the
algebra generated by $P_V$ and $\mathcal D_V$ is the full matrix algebra on the retained auxiliary
space. Equivalently, the blocked admissible transfer tensor is injective. Therefore the
corresponding blocked parent Hamiltonians satisfy a system-size-independent gap bound
\[
\inf_n \operatorname{gap}(H_{\mathrm{par}}(V_n))\ge m_0>0.
\]
Consequently no separate uniform parent-gap assumption remains in
\Cref{thm:gap-stability-thermo-limit}.
\end{theorem}

\begin{proof}
By \Cref{prop:strict-positive-c6-green}, every local admissible $C_6$ kernel has strictly positive
matrix entries. After finite blocking this yields a blocked operator $P_V$ whose matrix entries are
strictly positive on the retained basis. By
\Cref{lem:diagonal-weights-separate-basis}, the diagonal family $\mathcal D_V$ separates the
retained basis words. Hence every subspace invariant under the algebra
\[
\mathrm{Alg}(P_V,\mathcal D_V)
\]
must first be invariant under all separating diagonal matrices, and is therefore a coordinate
subspace spanned by a subset of basis vectors. But a strictly positive matrix sends every basis
vector to a vector with nonzero component in every retained basis direction, so no proper
coordinate subspace can remain invariant under $P_V$. Thus the generated algebra acts irreducibly.

Burnside's theorem now gives
\[
\mathrm{Alg}(P_V,\mathcal D_V)=M_n(\mathbb C)
\]
on the finite blocked auxiliary space. This is exactly the injectivity criterion after finite
blocking for the corresponding MPS / finitely-correlated transfer tensor. The standard
Fannes--Nachtergaele--Werner parent-Hamiltonian theorem therefore yields a uniform spectral gap
independent of system size. Hence
\[
\inf_n \operatorname{gap}(H_{\mathrm{par}}(V_n))>0.
\]
Then \Cref{thm:eta-zero-parent-domination} gives
\[
\operatorname{gap}(H_{\mathrm{adm}}(V_n))
\ge
\operatorname{gap}(H_{\mathrm{par}}(V_n))
\ge
m_0>0,
\]
and the thermodynamic-limit argument of \Cref{thm:thermodynamic-limit} yields a unique gapped vacuum
with exponential clustering.
\end{proof}

\begin{figure}[t]
\centering
\scriptsize
\begin{tikzpicture}[
    >=Latex,
    stagebox/.style={draw=blue!60!black, fill=blue!6, rounded corners, text width=2.2cm, minimum height=0.9cm, align=center},
    middlebox/.style={draw=orange!75!black, fill=orange!8, rounded corners, text width=2.2cm, minimum height=0.9cm, align=center},
    endbox/.style={draw=green!50!black, fill=green!8, rounded corners, text width=2.2cm, minimum height=0.9cm, align=center}
]
\node[stagebox] (pos) at (-5.4,0) {strictly positive\\local transport};
\node[stagebox] (sep) at (-2.9,0) {separating diagonal\\weights};
\node[middlebox] (irr) at (-0.4,0) {irreducibility\\of transfer algebra};
\node[middlebox] (burn) at (2.1,0) {Burnside\\$\Rightarrow M_n(\mathbb C)$};
\node[endbox] (inj) at (4.6,0) {injective blocked\\transfer tensor};
\node[endbox] (gap) at (7.1,0) {uniform parent gap\\and mass-gap stability};

\draw[->, thick] (pos) -- (sep);
\draw[->, thick] (sep) -- (irr);
\draw[->, thick] (irr) -- (burn);
\draw[->, thick] (burn) -- (inj);
\draw[->, thick] (inj) -- (gap);
\end{tikzpicture}
\caption{Analytic gap chain after the hardening step. The blocked transfer no longer assumes
injectivity a priori: strict positivity together with separating diagonal admissibility weights
forces irreducibility, hence full matrix algebra, and therefore injectivity and a uniform parent
gap.}
\label{fig:gap-hardening-chain}
\end{figure}
